\section{本章小结}
\label{sec:mfd-summary}

本章开头的问题是：经纬度在两极失效、角度绕圆一周会跳跃，这样的对象应当怎样记录？回答是放弃单张全局坐标，改用一族局部坐标图，并要求重叠处的转移映射光滑可逆。圆周的转移映射是平移，球极投影的转移映射是反演$\bm u/\lVert\bm u\rVert_2^2$，环面的转移映射是圆周转移映射的并排。

流形因此是一个整体空间：Hausdorff且第二可数，每点附近同胚于欧氏开集。光滑结构由光滑相容的坐标图决定，同一拓扑空间上可以有不同的光滑结构，例如$\mathbb R$上$x\mapsto x^3$给出的图与恒等图不相容。允许局部模型为半空间，就得到带边界流形；边界点的身份与图无关，边界本身是低一维的无边界流形，带角点的集合不属于这种光滑边界情形。

局部欧氏结构使流形的整体拓扑比一般拓扑空间规整。坐标球路径连通，所以路径分支都是既开又闭的，连通与路径连通等价；第二可数性又使连通分支至多可数。紧流形可以用有限张图覆盖，但正维数的紧无边界流形不能只用一张图，球面上的经纬度困难正源于此。

本章区分了三个层次的对象，表~\ref{tab:mfd-summary-records}用球面上的点$p=(0,\tfrac35,\tfrac45)$把它们并排列出。点是空间的元素，坐标的加法依赖图，不是点的运算。切向量是曲线速度的等价类，分量按转移映射的Jacobian矩阵变换；边界点处切空间用坐标记录定义，单侧可行的方向只占其中的半空间。余切向量是切向量的线性读数，系数按Jacobian矩阵的逆转置变换，在选定度量之前不能与切向量等同。

\begin{BookInlineTable}
  \centering
  \begin{tabular}{llll}
    \toprule
    对象 & 北极图 & 南极图 & 换算\\
    \midrule
    点$p$ & $(0,3)$ & $(0,\tfrac13)$ & $\tau$\\
    切向量$\partial_{u_1}|_p$ & $(1,0)$ & $(\tfrac19,0)$ & $\bm J$\\
    余切向量$\mathrm dh_p$ & $(0,0.12)$ & $(0,-1.08)$ & $\bm J^{-\top}$\\
    \bottomrule
  \end{tabular}
  \caption{球面上同一组对象在两张图中的记录，其中$\bm J=\operatorname{diag}(\tfrac19,-\tfrac19)$。}
  \label{tab:mfd-summary-records}
\end{BookInlineTable}

随点变化的对象在坐标中由分量记录，分量光滑即对象光滑。光滑函数的微分$\mathrm df_p$是余切向量，内点处的局部极值点微分为零。向量场作用于函数得到方向导数，两个向量场先后作用的交换差是一阶算子，给出Lie括号，它不依赖度量或联络。同一点的场值可以运算，不同点处的场值属于不同的切空间，坐标分量之差依赖图，不能当作它们的比较。光滑单位分解把各图上的构造加权拼成整体，依赖流形定义中的全部条件。

这些都是单个流形内部的结构。点从一个流形映到另一个流形时，切向量与余切向量怎样随之改变，是第~\ref{chap:manifold-maps}章的问题；比较不同点处的方向、度量长度和距离，要在第~\ref{chap:riemannian-geometry}章引入黎曼度量与联络。
